% concatenated from paper_arxiv.tex
\documentclass[11pt,oneside]{amsart}
%% ---------------------------------------------------------------
%% Preamble of mystyle2.sty, inlined for a self-contained submission.
%% ---------------------------------------------------------------
\usepackage{amsmath,amssymb,amsthm,mathtools}
\usepackage{bm}
\usepackage{enumitem}
\usepackage{xspace}
\usepackage[dvipsnames]{xcolor}
\usepackage[numbers,square,sort&compress]{natbib}
\usepackage{hyperref}
\hypersetup{
  colorlinks=true,
  linkcolor=MidnightBlue,
  citecolor=MidnightBlue,
  urlcolor=MidnightBlue
}

\emergencystretch=2em

\theoremstyle{plain}
\newtheorem{theorem}{Theorem}[section]
\newtheorem{proposition}[theorem]{Proposition}
\newtheorem{corollary}[theorem]{Corollary}
\newtheorem{lemma}[theorem]{Lemma}
\newtheorem{conjecture}[theorem]{Conjecture}

\theoremstyle{definition}
\newtheorem{definition}[theorem]{Definition}
\newtheorem{assumption}[theorem]{Assumption}
\newtheorem{example}[theorem]{Example}

\theoremstyle{remark}
\newtheorem{remark}[theorem]{Remark}

\newtheorem*{theorem*}{Theorem}

\usepackage[nameinlink,capitalize,noabbrev]{cleveref}
\crefname{theorem}{Theorem}{Theorems}
\Crefname{theorem}{Theorem}{Theorems}
\crefname{proposition}{Proposition}{Propositions}
\Crefname{proposition}{Proposition}{Propositions}
\crefname{corollary}{Corollary}{Corollaries}
\Crefname{corollary}{Corollary}{Corollaries}
\crefname{lemma}{Lemma}{Lemmas}
\Crefname{lemma}{Lemma}{Lemmas}
\crefname{definition}{Definition}{Definitions}
\Crefname{definition}{Definition}{Definitions}
\crefname{assumption}{Assumption}{Assumptions}
\Crefname{assumption}{Assumption}{Assumptions}
\crefname{example}{Example}{Examples}
\Crefname{example}{Example}{Examples}
\crefname{remark}{Remark}{Remarks}
\Crefname{remark}{Remark}{Remarks}

\newcommand{\NN}{\mathbb{N}}
\newcommand{\RR}{\mathbb{R}}
\newcommand{\CC}{\mathbb{C}}
\newcommand{\KK}{\mathbb{K}}
\newcommand{\TT}{\mathbb{T}}
\newcommand{\HU}{H^1(\Omega)}
\newcommand{\Hhalf}{H^{1/2}(\partial\Omega)}
\newcommand{\Hmhalf}{H^{-1/2}(\partial\Omega)}
\newcommand{\LDd}{L^2_\diamond(\partial\Omega)}
\newcommand{\iid}{\overset{\mathrm{i.i.d.}}{\sim}}
\newcommand{\norm}[1]{\left\lVert #1\right\rVert}
\newcommand{\abs}[1]{\left\lvert #1\right\rvert}
\newcommand{\calH}{\mathcal{H}}
\newcommand{\calK}{\mathcal{K}}
\newcommand{\calL}{\mathcal{L}}
\newcommand{\calM}{\mathcal{M}}
\newcommand{\calS}{\mathcal{S}}
\newcommand{\calT}{\mathcal{T}}
\newcommand{\calX}{\mathcal{X}}
\newcommand{\calY}{\mathcal{Y}}
\newcommand{\calZ}{\mathcal{Z}}
\newcommand{\eg}{e.g.\xspace}
\newcommand{\ie}{i.e.\xspace}

\DeclareMathOperator{\diag}{diag}
\DeclareMathOperator{\dist}{dist}
\DeclareMathOperator{\ran}{ran}
\DeclareMathOperator{\supp}{supp}
\DeclareMathOperator{\vol}{vol}
\DeclareMathOperator*{\essinf}{ess\,inf}
\DeclareMathOperator*{\esssup}{ess\,sup}

\newcommand{\restr}[2]{\left.#1\right|_{#2}}
%% --------------------- end of inlined preamble -------------------

\usepackage{geometry}

\title[Finite random measurements]{Analytic inverse problems with\\ finitely many random measurements}

\author{Giovanni S.~Alberti}
\address{DIMA, Department of Mathematics, University of Genoa, Via Dodecaneso 35, 16146 Genova, Italy}
\email{giovanni.alberti@unige.it}

\author{Damiano Poletti}
\address{DIMA, Department of Mathematics, University of Genoa, Via Dodecaneso 35, 16146 Genova, Italy}
\email{damiano.poletti@unige.it}

\author{Simone Sanna}
\address{DIMA, Department of Mathematics, University of Genoa, Via Dodecaneso 35, 16146 Genova, Italy}
\email{simone.sanna@edu.unige.it}

\author{Matteo Santacesaria}
\address{DIMA, Department of Mathematics, University of Genoa, Via Dodecaneso 35, 16146 Genova, Italy}
\email{matteo.santacesaria@unige.it}
\makeatletter
\def\shortauthors{Alberti et al.}
\makeatother

\subjclass[2020]{35R30, 14P15, 60B11}
\keywords{Nonlinear inverse problems, random measurements, sample complexity, analytic maps, Calder\'on problem, inverse scattering}
\hypersetup{
  pdftitle={Analytic inverse problems with finitely many random measurements},
  pdfauthor={Giovanni S. Alberti, Damiano Poletti, Simone Sanna, and Matteo Santacesaria}
}

\begin{document}

\begin{abstract}
While infinite-dimensional inverse problems are traditionally analyzed assuming continuous data, practical applications rely on finitely many discrete measurements. Recent deterministic approaches establish that unknowns belonging to a $d$-dimensional model class can be stably recovered from finitely many measurements. However, for severely ill-posed problems, such as the Calderón problem and inverse scattering, the known constructions may require a number of measurements that is exponential in $d$. We show that random sampling reduces this count dramatically if one asks only for exact identifiability. By exploiting the analytic geometry of the forward maps, we prove that, whenever the infinite-data problem is injective on the model class, $2d+1$ random scalar measurements determine the unknown uniquely, almost surely. Applications are given to the Calderón problem, with both infinite- and finite-dimensional boundary sampling, and to inverse medium scattering from randomly sampled far-field values.
\end{abstract}

\maketitle

\section{Introduction}\label{sec:introduction}

\subsection{Finite data in inverse problems}
An inverse problem seeks to recover an unknown parameter from indirect observations, often through a partial differential equation or, more generally, through a nonlinear map between infinite-dimensional spaces \cite{Isakov2017-qr,kirsch2011introduction}. The continuous model is natural for the analysis, but practical data acquisition produces only finitely many scalar observations. This leads to a basic sample-complexity question:
\begin{center}
\emph{How many scalar measurements are sufficient to identify every admissible parameter?}
\end{center}

We address this question under a finite-dimensional geometric prior. The admissible parameters are assumed to lie in a $d$-dimensional embedded real-analytic manifold inside an ambient Banach space. Under this geometric prior, we derive sufficient bounds on the number of measurements needed in order to guarantee exact reconstruction of an arbitrary unknown in the parameter space. 

\subsection{Main results at a glance}\label{sec:main-results-glance}
Recently, it has been shown that finite deterministic measurements can guarantee stable recovery for finite-dimensional unknowns \cite{alberti2022infinite}. Nevertheless, for notoriously ill-posed problems such as the Calder\'on problem \cite{alessandrini1988stable,mandache2001exponential} or inverse scattering \cite{stefanov1990stability}, this deterministic sampling approach may require an exponential number of measurements relative to the dimension of the parameter space, see \cite[Example 4]{alberti2022infinite}.

Our primary contribution is demonstrating that a probabilistic approach to sampling the measurements drastically reduces the required sample complexity for exact recovery. We establish that exact recovery can be achieved almost surely with a number of measurements that is merely linear in the dimension of the parameter space.

In order to highlight the impact of this approach, we first describe the two principal applications of our main result \cref{thm:main} informally. Together with them we also illustrate the recovery result for sparse unknowns \cref{thm:sparse}.
\vspace{.1cm}
\paragraph{\emph{The Calder\'on problem.}}
Let $\Omega\subset\RR^m$, $m\geq 2$, be a bounded Lipschitz domain and let $\mathcal M\subset L^\infty(\Omega)$ be a given $d$-dimensional analytic model class of uniformly positive conductivities. For  $\sigma\in\mathcal M$, the Neumann-to-Dirichlet map $\Lambda_\sigma$ sends a mean-free boundary current $f$ to the corresponding mean-free boundary voltage, see \eqref{eq:ND-map}. Given independent Gaussian boundary functions $f_i,g_i$, $i=1,\dots,M$, we observe only the scalar matrix elements
\begin{equation}\label{eq:intro-calderon-measurements}
    \bigl(\langle \Lambda_\sigma f_i,g_i\rangle_{L^2(\partial\Omega)}\bigr)_{i=1}^M.
\end{equation}
We prove that, if the full forward map $\sigma\mapsto\Lambda_\sigma$ is injective on the fixed model class, then any $M\geq 2d+1$ independently sampled observations of this form identify every conductivity in the class almost surely; see \cref{thm:calderon-gaussian}.

We also prove a finite-dimensional boundary-sampling result. On a compact parameter set, injectivity of the full forward map and of its differential implies that a sufficiently large finite-dimensional truncation of the boundary space preserves both properties. Random currents and voltages drawn from that fixed finite-dimensional space then yield the same $2d+1$ sample bound; see \cref{thm:calderon-finite}.
\vspace{.1cm}
\paragraph{\emph{Inverse scattering.}}
Let $n$ be the sought-after refractive index of an inhomogeneous medium supported within a ball $B\subset\RR^3$. For an incident direction $\theta\in S^2$, we probe the medium using the incident plane wave
\begin{equation}\label{eq:plane-wave}
    u_\theta^i(x)=e^{ikx\cdot\theta}.
\end{equation}
As this wave interacts with the inhomogeneous medium, it produces a scattered wave. The sum of the incident wave and this scattered one constitutes the total field, denoted by $u_{n,\theta}$. Then, detectors are placed far away to record the scattered energy. These measurements are captured by the far-field pattern $u_n^\infty(\widehat{x},\theta)$, which represents the scattering amplitude in an observation direction $\widehat{x}\in S^2$. It is governed by the integral equation \eqref{eq:far-field}. The corresponding inverse problem consists of recovering $n$ after probing the medium with incident waves in every possible direction $\theta\in S^2$ and measuring the corresponding far-field pattern on the whole $S^2$. Assuming that the full far-field map $n\mapsto u^\infty_n$  is injective on a given $d$-dimensional analytic model class, then $2d+1$ independent random pairs $(\widehat x_i,\theta_i)$ identify every admissible refractive index almost surely through the measurements $u_n^\infty(\widehat x_i,\theta_i)$; see \cref{thm:scattering}.
\vspace{.1cm}
\paragraph{\emph{Sparse recovery.}} In many applications, the parameters we seek to recover are sparse with respect to a fixed basis. The abstract result \cref{thm:main} allows us to derive a recovery guarantee for $s$-sparse parameters where the number of measurements depends linearly on the sparsity level $s$ that is typically much smaller than the dimension of the ambient space $d$. This can be interpreted as an exact, nonlinear analog to noise-free compressed sensing \cite{Foucart2013-od}. Indeed, classical compressed sensing in linear scenarios typically requires $\mathcal O(s\log(d/s))$ measurements for stable recovery \cite{candes2006robust,donoho2006compressed}. Our geometric framework guarantees only exact identifiability (without convex recovery) with $4s+1$ measurements (eliminating the logarithmic term), but the setting is fully nonlinear, see \cref{thm:sparse}.

\subsection{Geometric mechanism}
The effectiveness of this approach stems from a geometric intuition. In both the examples we presented above, the forward model is analytic. Therefore, the set of observations that do not distinguish between two given parameters is the zero-level set of an analytic function.

If the underlying infinite-dimensional problem is known to be uniquely solvable, this analytic function cannot be identically zero. Consequently, its zero set has positive codimension in the sampling space. To obtain a uniform statement over all parameter pairs, we form an incidence set over the $2d$-dimensional pair manifold and estimate the dimension of its projection onto the space of measurement systems. Once the number of scalar measurements exceeds $2d$, the uninformative measurement systems have positive codimension and hence probability zero.

This mechanism is reminiscent of Whitney-type and prevalent embedding theorems \cite{whitney1936}, but the observables in the present work are prescribed by the physics of the inverse problem rather than chosen from an unrestricted class.


\subsection{Related work}\label{sec:related-work}
Finite data uniqueness and stability for inverse problems have been studied in a number of deterministic settings. For the Calder\'on problem and related elliptic inverse problems, see \cite{alberti2019calderon,alberti2022calderon,harrach2019uniqueness,hanke2024lipschitz,ruland2019fractional,ruland2022runge}; for inverse scattering with strongly structured unknowns, see \cite{blaasten2020recovering,blaasten2021corners}. Fully discrete regularization in the Calder\'on and scattering settings is studied in \cite{felisi2024full,di2026discretization}. Generic finite measurements obtained through singularity-theoretic arguments also appear in the fractional Calder\'on problem with drift \cite{cekic2020calderon}. Abstract finite-measurement stability results for finite-dimensional subspaces and manifolds were developed in \cite{alberti2022infinite,alberti2022inverse}. The recent work \cite{carstea2026h} derives H\"older stability and qualitative finite determinacy from exact uniqueness in finite-dimensional analytic inverse problems. Our result is complementary: it gives an explicit, universal almost sure count of random scalar measurements, but does not establish stability. %

The geometric antecedents include Whitney's embedding theorem \cite{whitney1936}, Ma{\~n}{\'e}'s projection theorem \cite{mane1981dimension}, Takens' delay-coordinate theorem \cite{takens1981detecting}, and the ``embedology'' framework of Sauer, Yorke, and Casdagli \cite{sauer1991embedology}. Generic and prevalent finite-dimensional projections of compact or fractal subsets of infinite-dimensional spaces were studied in \cite{hunt1992prevalence,hunt1999regularity,robinson2009linear}. Random stable embeddings of smooth manifolds and more general structured sets are treated in \cite{baraniuk2009random,dirksen2016dimensionality}. Those works allow broad classes of observables or random linear projections; here the measurements belong to a fixed analytic family generated by an inverse problem forward map.

Our incidence set argument is also related to algebraic identifiability in phase retrieval and algebraic compressed sensing, where the dimension of a bad measurement variety controls generic injectivity \cite{conca2015algebraic,breiding2021algebraic}. In signal processing, random measurements of sparse sets, manifolds, and generative-model ranges lead to stable recovery under quantitative geometric assumptions \cite{candes2006robust,donoho2006compressed,Foucart2013-od,bora2017compressed}. These principles have also been applied to inverse scattering \cite{gilbert-levinson-schotland-2020,alberti-petit-santacesaria}. The present theorem instead concerns exact global injectivity for structured nonlinear measurements and uses real-analytic geometry rather than concentration inequalities.

Finally, randomization is increasingly used for data selection and sketching in PDE-constrained inverse problems. We refer to \cite{jin2024unique,HELLMUTH2026} and the references therein. Statistical and Bayesian inverse problems pursue a different objective: the unknown is assigned a prior distribution and uncertainty is quantified through a posterior law \cite{kaipio2005statistical,stuart2010inverse}; see \cite{abraham2020statistical} for the statistical Calder\'on problem.

\subsection{Open questions and comments}

The count $2d+1$ obtained above concerns scalar observations rather than
physical experiments. In the Calder\'on problem, one applied current $f$
produces the whole boundary voltage $\Lambda_\sigma f$, while in inverse
scattering one incident direction $\theta$ produces the full far-field pattern
$u_n^\infty(\cdot,\theta)$. It is therefore natural to ask whether these
vector-valued data permit recovery from fewer experiments, possibly from a
single generic input--response pair. In
\cite{boulle2026zero}, the authors prove a sharp zero--one law of this type for linearly
parameterized analytic systems: either no input identifies the coefficients,
or almost every input sampled from a nondegenerate Gaussian measure does.
Extending this dichotomy to nonlinear parameter-to-operator maps may reveal
if one-shot recovery in the Calder\'on and
scattering problems is possible.

The present results concern uniqueness only. The approach of
 \cite{carstea2026h}, based on {\L}ojasiewicz-type inequalities,
suggests that, for almost every fixed sampling system, analytic injectivity
on a compact model class may yield a sample-dependent H\"older estimate. A
quantitative version should, for every $\varepsilon>0$, provide deterministic
constants $C_\varepsilon>0$ and exponents
$\alpha_\varepsilon\in(0,1]$ for which a uniform H\"older estimate holds with
probability at least $1-\varepsilon$. The stability constant should be
expected to deteriorate as $\varepsilon\downarrow0$. Lipschitz stability,
corresponding to $\alpha_\varepsilon=1$, may be possible under stronger
uniform lower bounds for the differential of the sampled forward map.

Reconstruction is another open direction. In the sparse setting of
\cref{thm:sparse}, one could adapt nonlinear restricted-isometry and
iterative hard-thresholding techniques \cite{blumensath2013compressed}, as
well as sparsity-promoting, support-informed proximal methods for nonlinear
EIT \cite{lazzaro2024oracle}. The main challenge is to derive robust
geometric conditions for the present random measurements that imply global
convergence, possibly from arbitrary initializations, and stability with
respect to noise.

Finally, it would be interesting to extend the incidence-set argument to
other nonlinear inverse problems, especially those governed by semilinear
or quasilinear PDEs. The key requirements are analytic dependence on the
unknown and the experimental input, together with a suitable full-data
separation property.

\subsection{Organization of the paper}
The main abstract theorem, \cref{thm:main}, treats two measurement architectures. The first consists of a jointly analytic scalar measurement map indexed by a finite-dimensional analytic manifold. The second consists of multilinear measurements generated by a nondegenerate Gaussian measure on a Hilbert space. We then prove a sparse variant, \cref{thm:sparse}. The applications to the Calder\'on problem and inverse scattering are stated in \cref{sec:applications}. The geometric proof of the abstract results is given in \cref{sec:proof-main}, the application to the Calder\'on problem in \cref{sec:proof-calderon}, and the application to inverse scattering  in \cref{sec:proof-scattering}.

\section{Random analytic measurements}\label{sec:framework}

Throughout the paper, $X$ is a real Banach space with norm $\|\cdot\|_X$, $A\subset X$ is open, and $\calM\subset X$ is a second-countable\footnote{Second countability is a genuine restriction and not a consequence of the
ambient structure: the spaces $X$ arising in the applications, such as $L^\infty(\Omega)$,
are not separable. It is used only through \cref{lem:semianalytic-exhaustion}.} embedded real-analytic manifold of finite dimension $d$. Complex Banach spaces are always regarded as real Banach spaces when analyticity with respect to the parameter is discussed. The scalar field is denoted by $\KK\in\{\RR,\CC\}$.


We first make precise the two notions of analyticity used throughout.

\begin{definition}\label[definition]{def:analytic}
Let $Y_1,Y_2$ be real Banach spaces and let $U\subset Y_1$ be open. A map
$G\colon U\to Y_2$ is \emph{real analytic} if for every $y\in U$ there exist $\rho>0$ and
continuous symmetric $j$-linear maps $G_j\colon Y_1^j\to Y_2$, $j\geq0$, with
$\sum_{j\geq0}\norm{G_j}\rho^j<\infty$ and
\[
    G(y+h)=\sum_{j\geq0}G_j(h,\ldots,h),
    \qquad \norm{h}_{Y_1}<\rho .
\]
A subset $\calM\subset X$ is an \emph{embedded real-analytic manifold of dimension $d$} if
for every $a\in\calM$ there are an open set $V\subset\RR^d$ and a real-analytic map
$\varphi\colon V\to X$ such that $d\varphi_t\colon\RR^d\to X$ is injective for every
$t\in V$ and $\varphi$ is a homeomorphism from $V$ onto an open neighbourhood of $a$ in
$\calM$, the latter being endowed with the topology induced by $X$. We call such a
$\varphi$ a \emph{chart} of $\calM$ and we set $T_a\calM=d\varphi_{\varphi^{-1}(a)}(\RR^d)$.
\end{definition}

The transition map between two charts is real analytic. Indeed, if $\varphi,\psi$ are
charts around $a$ and $s=\psi^{-1}(a)$, then $d\psi_s$ has finite-dimensional, hence
complemented, range, so there is a bounded linear map $L\colon X\to\RR^d$ for which
$L\circ d\psi_s$ is invertible; the analytic inverse function theorem applied to
$L\circ\psi$ shows that $\psi^{-1}=(L\circ\psi)^{-1}\circ L$ is real analytic near $a$,
and therefore so is $\psi^{-1}\circ\varphi$. Thus $\calM$ carries a well-defined
$d$-dimensional real-analytic structure, and the inclusion $\calM\hookrightarrow X$ is
real analytic in every chart.


We consider two sampling architectures.

\subsection{Finite-dimensional analytic sampling}\label{sec:finite-analytic-sampling}
Let $Z$ be a connected real-analytic manifold of finite dimension $q\geq1$, endowed with a smooth positive volume measure $\vol_Z$. Let $\mu$ be a Borel probability measure such that
\begin{equation}\label{eq:mu-ac}
    \mu\ll \vol_Z.
\end{equation}
Let
\begin{equation}\label{eq:analytic-measurement-map}
    \mathfrak m\colon A\times Z\longrightarrow\KK
\end{equation}
be a real-analytic measurement map. We impose that any parameter $a\in A\cap\calM$ is uniquely determined when the measurements associated with every possible sample in $Z$ are observed. In other words, we assume 
\begin{equation}\label{eq:analytic-separation}
    \mathfrak m(a_1,\cdot)-\mathfrak m(a_2,\cdot)\not\equiv0\ \text{on }Z,\quad\text{for every}~a_1\neq a_2\ \text{in }A\cap\calM.
\end{equation}

This formulation covers pointwise sampling of analytic data, as in inverse scattering, and scalar matrix elements restricted to a finite-dimensional input space, as in \cref{thm:calderon-finite}.

\subsection{Gaussian multilinear sampling}\label{sec:gaussian-sampling}
Let $H$ be a separable Hilbert space over $\KK$, let $Z$ be an infinite-dimensional separable real Hilbert space, and let $r\geq1$. We write $\calL^r(H;\KK)$ for the Banach space of bounded $\KK$-multilinear forms on $H^r$. Let the forward map
\begin{equation}\label{eq:multilinear-forward}
    F\colon A\longrightarrow\calL^r(H;\KK)
\end{equation}
be real-analytic and injective on $A\cap\calM$. Let $\psi\colon Z\to H$ be bounded and real linear, and assume that
\begin{equation}\label{eq:dense-span}
    \overline{\operatorname{span}_{\KK}\psi(Z)}=H.
\end{equation}
The real-analytic scalar measurement map is
\begin{equation}\label{eq:gaussian-measurement-map}
    \mathfrak m\colon A\times Z^r\to \KK,\qquad \mathfrak m(a,z)
    =F(a)\bigl(\psi z^1,\ldots,\psi z^r\bigr).
\end{equation}
Since $\psi$ is only real linear, $\mathfrak m(a,\cdot)$ is $\RR$-multilinear on $Z^r$ even when $\KK=\CC$; this is all that will be used.
We abuse notation and we denote the measurement map by $\mathfrak m$ both for the finite-dimensional analytic sampling (eq.\ \eqref{eq:analytic-measurement-map}) and for the Gaussian multilinear sampling (eq.\ \eqref{eq:gaussian-measurement-map}).

Let $\gamma$ be a Gaussian Borel probability measure on $Z$. We call $\gamma$ nondegenerate if its covariance operator has trivial kernel. Equivalently, every nonzero continuous linear functional on $Z$ has a nondegenerate one-dimensional Gaussian distribution under $\gamma$ \cite{Bogachev1998-oc,da2006introduction}. One sample is an $r$-tuple $z=(z^1,\ldots,z^r)\in Z^r$, distributed according to $\gamma^r$.
 

The injectivity of $F$ on $A\cap\calM$ and \eqref{eq:dense-span} imply the analogue of the point-separation condition \eqref{eq:analytic-separation}, namely
\begin{equation}\label{eq:gaussian-separation}
    \mathfrak m(a_1,\cdot)-\mathfrak m(a_2,\cdot)\not\equiv0\ \text{on }Z^r,\qquad\text{for every}~a_1\neq a_2\ \text{in }A\cap\calM.
\end{equation}

Indeed, assume that $\mathfrak m(a_1,z)=\mathfrak m(a_2,z)$ for every $z\in Z^r$ and set
$T=F(a_1)-F(a_2)\in\calL^r(H;\KK)$. By \eqref{eq:gaussian-measurement-map}, $T$ vanishes on
$\psi(Z)^r$; by $\KK$-multilinearity it vanishes on
$\bigl(\operatorname{span}_\KK\psi(Z)\bigr)^r$, and by continuity and \eqref{eq:dense-span}
it vanishes on $H^r$. Hence $F(a_1)=F(a_2)$, and the injectivity of $F$ on $A\cap\calM$
gives $a_1=a_2$.


\subsection{Uniform recovery}\label{sec:uniform-recovery}
For either architecture, denote the sample space and probability measure by $(\calZ,\nu)$: in the finite-dimensional case, $(\calZ,\nu)=(Z,\mu)$ and the measurement is \eqref{eq:analytic-measurement-map}; in the Gaussian case, $(\calZ,\nu)=(Z^r,\gamma^r)$ and the measurement is \eqref{eq:gaussian-measurement-map}.

We now state the main result of the work ensuring uniform identifiability of parameters lying within a $d$-dimensional real-analytic manifold using $2d+1$ scalar measurements arising from independent random samples on $(\calZ,\nu)$.

\begin{theorem}\label{thm:main}
Let $X$ be a real Banach space, $A\subset X$ be open, and $\calM\subset X$ be a second-countable embedded real-analytic manifold of finite dimension $d$. Assume one of the two sampling architectures described in \cref{sec:finite-analytic-sampling,sec:gaussian-sampling}, denote by $\mathfrak m$ the corresponding real-analytic scalar measurement map (see \eqref{eq:analytic-measurement-map} and \eqref{eq:gaussian-measurement-map}) satisfying the corresponding separation and nondegeneracy hypotheses (see \eqref{eq:analytic-separation} and \eqref{eq:gaussian-separation}). Let $M\geq 2d+1$, and let $z_1,\ldots,z_M\in\calZ$ be independent samples with law $\nu$. Then, with probability one, for every $a_1,a_2\in A\cap\calM$, 
\[
   \text{if } \mathfrak m(a_1,z_i)=\mathfrak m(a_2,z_i),\; \text{for } i=1,\dots,M, \quad \text{ then } a_1=a_2.
\]
\end{theorem}

Equivalently, the set of measurement systems that fail to distinguish at least one pair of distinct parameters has $\nu^M$-measure zero. 

The same arguments can be used to show non-uniform recovery, i.e.\ pointwise recovery, using $d+1$ random samples.

\begin{corollary}\label[corollary]{cor:nonuniform}
Under the hypotheses of \cref{thm:main}, fix $a^\dagger\in A\cap\calM$. If $M\geq d+1$, then, with probability one, for every $a\in A\cap\calM$, 
\[ \text{if } 
    \mathfrak m(a,z_i)=\mathfrak m(a^\dagger,z_i),
    \; \text{for } i=1,\ldots,M, \quad \text{ then } a=a^\dagger.
\]
\end{corollary}

\begin{remark}\label[remark]{rem:complex-measurements}
All dimensions in the proof are real dimensions. A nontrivial complex-valued real-analytic equation has a zero set of real codimension at least one, which is all that is used here. Thus, the bound $2d+1$ is a sufficient bound for $\CC$-valued measurements and is not asserted to be optimal.
\end{remark}

\subsection{Sparse parameters}\label{sec:sparse}
Together with the recovery guarantee for generic finite-dimensional parameters stated above, we also provide an application of the same techniques for the identifiability of sparse unknowns.

The focus on recovering sparse unknowns is motivated by the inherent structure of most natural signals. Consequently, these signals typically admit sparse representations \cite{zhang2015survey} when expressed in appropriate bases, such as wavelets \cite{mallat1999wavelet} or shearlets \cite{kutyniok2012introduction}. Mathematically, this implies that high-dimensional signals often possess a much lower intrinsic dimension. Following the established paradigm of compressed sensing \cite{Foucart2013-od}, this allows for signal recovery using significantly fewer measurements than the ambient dimension dictates. Such a reduction in sample complexity is crucial in contexts where data acquisition is slow or expensive.

In mathematical terms, we restrict $\mathcal M\subset X$ to be a $d$-dimensional real linear subspace and we deal with unknown parameters $a$ that are $s$-sparse with respect to a given basis $(e_j)_{j=1}^d$ of $\calM$, namely
\[
    a = \sum_{j=1}^d \alpha_j e_j,
\]
with $|\operatorname{supp}(\alpha)|\le s$ where 
\[
    \operatorname{supp}(\alpha)=\{j:\alpha_j\neq0\},\quad\alpha\in\RR^d.
\]
For $1\leq s\leq d$, define
\begin{equation}\label{eq:sparse-set}
    \calM_s
    =\bigcup_{\substack{I\subset\{1,\ldots,d\}\\ |I|\leq s}}
      \operatorname{span}\{e_j:j\in I\}
\end{equation}
as the union of every $s$-dimensional subspace of $\calM$ generated by $s$ basis elements among $(e_j)_{j=1}^d$. 

Relying on the abstract result \cref{thm:main}, we derive a uniform recovery result for $s$-sparse unknowns using $4s+1$ random measurements. Notice that in many practical scenarios $s$ is much smaller than $d$ yielding $2s<d$, hence the sample complexity for the sparse case typically is significantly better than the general one for $d$-dimensional parameters.
\begin{theorem}\label{thm:sparse}
Let $X$ be a real Banach space, $A\subset X$ be open, $\calM\subset X$ be a real linear subspace of finite dimension $d$, and $(e_j)_{j=1}^d$ be a basis for $\calM$. Assume one of the two sampling architectures described in \cref{sec:finite-analytic-sampling,sec:gaussian-sampling}, denote by $\mathfrak m$ the corresponding real-analytic scalar measurement map (see \eqref{eq:analytic-measurement-map} and \eqref{eq:gaussian-measurement-map}) satisfying the corresponding separation and nondegeneracy hypotheses (see \eqref{eq:analytic-separation} and \eqref{eq:gaussian-separation}). Let
\begin{equation}\label{eq:sparse-bound}
    M\geq 2\min\{d,2s\}+1,
\end{equation}
and let $z_1,\dots,z_M\in\calZ$ be independent samples with law $\nu$. Then, with probability one, for every $a_1,a_2\in A\cap\calM_s$,
\[ \text{if }
   \mathfrak m(a_1,z_i)=\mathfrak m(a_2,z_i),\; \text{for } i=1,\dots,M, \quad \text{then } a_1=a_2.
\]
\end{theorem}

The theorem is an exact, noiseless statement. Unlike quantitative compressed sensing, it does not give robustness to noise or an algorithmic recovery guarantee. Its advantage is that the measurements may be nonlinear and constrained by an inverse problem forward map.

\section{Applications}\label{sec:applications}

\subsection{The Calder\'on problem}\label{sec:calderon}
The Calder\'on problem \cite{calderon2006inverse} is the mathematical foundation of electrical impedance tomography (EIT) \cite{cheney1999electrical,borcea2002electrical}. It asks whether the electrical conductivity of a body can be uniquely determined by performing current-to-voltage measurements on the boundary.

Let $\Omega\subset\RR^m$, $m\geq2$, be a bounded connected Lipschitz domain. All function spaces in this subsection are real. Set
\[
    L^\infty_+(\Omega)
    =\{\sigma\in L^\infty(\Omega):\essinf_\Omega\sigma>0\}.
\]
For $\sigma\in L^\infty_+(\Omega)$ and $f\in L^2_\diamond(\partial\Omega)$, let $u_{\sigma,f}\in H^1(\Omega)$ be the weak solution, normalized by $\int_{\partial\Omega}u_{\sigma,f}\,ds=0$, of
\begin{equation}\label{eq:conductivity-equation}
  \left\{\begin{aligned}
    &-\nabla\cdot(\sigma\nabla u)=0&&\mathrm{in}~\Omega,\\
    &\sigma\partial_\nu u=f&&\mathrm{on}~\partial\Omega.
    \end{aligned}\right.  
\end{equation}
Here
\[
    L^2_\diamond(\partial\Omega)
    =\left\{f\in L^2(\partial\Omega):\int_{\partial\Omega}f\,ds=0\right\}.
\]
The boundary measurements are modeled via the Neumann-to-Dirichlet map
\begin{equation}\label{eq:ND-map}
    \Lambda_\sigma\colon L^2_\diamond(\partial\Omega)\to L^2_\diamond(\partial\Omega),\qquad \Lambda_\sigma f=\restr{u_{\sigma,f}}{\partial\Omega}.
\end{equation}
The Calder\'on problem involves recovering $\sigma$ from $\Lambda_\sigma$.

It is well established that the problem is uniquely solvable for $m=2$ \cite{nachman1996global, astala2006calderon} (see also \cite{seo1995uniqueness} for a uniqueness result in the planar case with finitely many measurements) or for sufficiently smooth conductivities \cite{sylvester1987global, haberman2013uniqueness,haberman2015uniqueness,caro2016global}. It has been shown that it is only logarithmically stable \cite{alessandrini1988stable,mandache2001exponential}. Lipschitz type stability estimates have been shown for finite-dimensional unknowns with infinitely many measurements \cite{alessandrini2005lipschitz, rondi2006remark,bacchelli2006lipschitz,beretta2011lipschitz,aspri2022lipschitz} and with finitely many measurements \cite{harrach2019uniqueness,alberti2022infinite,alberti2022inverse,hanke2024lipschitz}. 

The map $\sigma\mapsto\Lambda_\sigma$ is real-analytic, see \cref{prop:nd-analytic}. Hence, assuming the injectivity of the forward model with full-data, we can apply the abstract recovery result, \cref{thm:main}, in the setting of \cref{sec:gaussian-sampling}, with:
\begin{itemize}
    \item $X=L^\infty(\Omega)$,
    \item $A=L^\infty_+(\Omega)$,
    \item $Z=H=L^2_\diamond(\partial\Omega)$,
    \item $r=2$,
    \item $F(\sigma)(f,g)=\langle\Lambda_{\sigma} f,g\rangle_{L^2(\partial\Omega)}$ for $\sigma\in A$ and $f,g\in H$,
    \item and $\psi=\operatorname{Id}\colon Z\to H$.
\end{itemize}
We obtain the following uniqueness guarantee using finitely many random scalar measurements with boundary samples from $L^2_\diamond(\partial\Omega)$.

\begin{theorem}\label{thm:calderon-gaussian}
Let $\calM\subset L^\infty(\Omega)$ be a second-countable embedded real-analytic manifold of dimension $d$, and assume that
\[
    \sigma\longmapsto\Lambda_\sigma
\]
is injective on $L^\infty_+(\Omega)\cap\calM$. Let $\gamma$ be a Gaussian measure on $L^2_\diamond(\partial\Omega)$ with injective covariance. Let $M\geq2d+1$ and let $(f_1,g_1),\dots,(f_M,g_M)\in L^2_\diamond(\partial\Omega)\times L^2_\diamond(\partial\Omega)$ be independent samples with law $\gamma\otimes\gamma$. Then, with probability one, for every $\sigma_1,\sigma_2\in L^\infty_+(\Omega)\cap\calM$, 
\[ \text{if }
   \langle\Lambda_{\sigma_1} f_i,g_i\rangle_{L^2(\partial\Omega)}=\langle\Lambda_{\sigma_2} f_i,g_i\rangle_{L^2(\partial\Omega)},\; \text{for } i=1,\dots,M,\quad \text{then } \sigma_1=\sigma_2.
\]
\end{theorem}

The inverse problem specific hypothesis is full-data injectivity on the chosen model class. As discussed above, such injectivity is known in dimension two for bounded scalar conductivities \cite{astala2006calderon} and, in higher dimensions, under several regularity or structural assumptions; see \cite{nachman1996global,sylvester1987global,alessandrini2005lipschitz,alessandrini2017lipschitz,harrach2019uniqueness,haberman2013uniqueness,haberman2015uniqueness,caro2016global}.

\subsubsection{Finite-dimensional boundary sampling}\label{sec:calderon-finite}
In the previous section we established a recovery guarantee for the Calder\'on problem where the input data are sampled from the infinite-dimensional space $L^2_\diamond(\partial\Omega)$.

However, since the parameter space is finite-dimensional, it is possible to appropriately truncate an orthonormal basis $(e_j)_{j\in\NN}$ for $L^2_\diamond(\partial\Omega)$ and consider the finite-dimensional space
\begin{equation}\label{eq HN}
    H_N=\operatorname{span}\{e_1,\dots,e_N\}\subset\LDd,\quad\text{for some}~N\in\NN,
\end{equation}
as the sampling space.

The following theorem relies on a uniform truncation statement, see \cref{subsection:uniform-preservation-injectivity-truncation}. The truncation must preserve not only injectivity on the compact search set, but also injectivity of the differential of the truncated forward map. In the following, we denote by $\calK\bigl(L^2_\diamond(\partial\Omega)\bigr)$ the set of compact operators acting between $L^2_\diamond(\partial\Omega)$ and itself.

\begin{theorem}\label{thm:calderon-finite}
Let $\calM\subset L^\infty(\Omega)$ be a second-countable embedded real-analytic manifold of dimension $d$, and let
\[
    K\subset L^\infty_+(\Omega)\cap\calM
\]
be compact. Assume that:
\begin{enumerate}[label=\textup{(\roman*)},leftmargin=*]
    \item $\sigma\mapsto\Lambda_\sigma$ is injective on $K$;
    \item for every $\sigma\in K$, the differential
    \[
        d\Lambda_\sigma\colon T_\sigma\calM\longrightarrow
        \calK\bigl(L^2_\diamond(\partial\Omega)\bigr)
    \]
    is injective.
\end{enumerate}
Then there exists $N_0\in\NN$ such that, for every $N\geq N_0$, the following holds. Let $\rho_N$ be a probability measure on $H_N$ that is absolutely continuous with respect to the Lebesgue measure. Let $M\geq2d+1$ and let $(f_1,g_1),\dots,(f_M,g_M)\in H_N\times H_N$ be independent samples with law $\rho_N\otimes\rho_N$. Then, with probability one, for every $\sigma_1,\sigma_2\in K$,
\[ \text{if }
   \langle\Lambda_{\sigma_1} f_i,g_i\rangle_{L^2(\partial\Omega)}=\langle\Lambda_{\sigma_2} f_i,g_i\rangle_{L^2(\partial\Omega)},\; \text{for } i=1,\dots,M,\quad \text{then } \sigma_1=\sigma_2.
\]
\end{theorem}

\begin{remark}\label[remark]{rem:compact-calderon-set}
A common choice is
\[
    K=\calM\cap\left\{\sigma\in L^\infty(\Omega):\lambda^{-1}\leq\sigma\leq\lambda\ \text{a.e.\ in}~\Omega\right\}
\]
for some $\lambda>1$, provided this intersection is compact in $\calM$. For instance, this is the case whenever $\calM\subset L^\infty(\Omega)$ is a finite-dimensional real linear subspace. However, compactness is an assumption; it does not follow merely from boundedness in the ambient $L^\infty$-norm for an arbitrary embedded manifold.
\end{remark}

\begin{remark}\label[remark]{rem:calderon-hypotheses}
Implementing this finite-dimensional truncation relies on the injectivity of both the forward operator, discussed earlier, and the linearized operator. The injectivity of the latter follows from the completeness of products of suitable solutions across several smooth settings \cite{sylvester1987global,paivarinta2003complex}; for structured conductivity classes, we refer to \cite{lechleiter2008newton,harrach2010exact}. The hypotheses of \cref{thm:calderon-finite} must be explicitly verified for the specific tangent spaces of the underlying model.

\end{remark}

\subsection{Inverse medium scattering}\label{sec:scattering}
The inverse medium scattering problem \cite{colton1998inverse, kirsch2011introduction} involves recovering an unknown complex-valued refractive index of a medium from far-field pattern measurements.

The mathematical model is the following, see \cite[Chapter 8]{colton1998inverse} for further details. Let $k>0$ be the fixed wavenumber and $n \in L^\infty(\RR^3;\CC)$ be a complex-valued refractive index. We assume the medium of interest is supported within a known bounded ball $B\subset\RR^3$, namely,  $n=1$ in $\RR^3\setminus B$. Let the open set of admissible parameters be
\begin{equation}\label{eq:scattering-admissible-open}
    \mathcal A_k
    =\{n\in L^\infty(B;\CC):I-V_n\text{ is invertible on }L^2(B)\},
\end{equation}
where $V_n\colon L^2(B)\to L^2(B)$ is the volume potential operator \cite[Theorem 8.2]{colton1998inverse} defined by
\begin{equation}\label{eq:volume-potential}
    V_n v
    =k^2\int_B G_k(\cdot,y)(n(y)-1)v(y)\,dy,
\end{equation}
where
\[
    G_k(x,y)=\frac{e^{ik|x-y|}}{4\pi|x-y|}
\]
is the outgoing fundamental solution of $\Delta+k^2$. For $n\in\mathcal A_k$ and for an incident plane wave $u^i_\theta(x)=e^{ik x\cdot \theta}$ that solves
\begin{equation}\label{eq helm}
    \Delta u_\theta^i +k^2 u_\theta^i=0\quad\text{in}~\RR^3,
\end{equation} 
we denote by $u_{n,\theta}$ the corresponding total field fulfilling 
\begin{equation}\label{eq phisical model scattering}
    \left\{\begin{aligned}
    &\Delta u_{n,\theta}+k^2nu_{n,\theta}=0&&\mathrm{in}~\RR^3,\\
    &u_{n,\theta}=u_\theta^i+u_{n,\theta}^s&&\mathrm{in}~\RR^3,
    \end{aligned}\right.   
\end{equation}
where $u_{n,\theta}^s$ is the scattered field satisfying the Sommerfeld radiation condition
\begin{equation}\label{eq sommerfeld cond}
    \lim_{r\to+\infty}r\bigg(\partial_r u^s_{n,\theta}-iku^s_{n,\theta}\bigg)=0,
\end{equation}
where $\partial_r$ denotes the radial derivative. Notice that,
by the Lippmann--Schwinger equation \cite[Equation 8.13]{colton1998inverse}, we have
\begin{equation}\label{eq:lippmann-schwinger}
    u^s_{n,\theta}=k^2\int_{B}G_k(\cdot,y)(n(y)-1)u_{n,\theta}(y)\,dy=V_n u_{n,\theta},
\end{equation}
and so, by the second equation in \eqref{eq phisical model scattering}, the total field is well-defined and uniquely given by
\[
    u_{n,\theta}=(I-V_n)^{-1}u_\theta^i.
\]
The corresponding far-field pattern on $S^2$ is given by \begin{equation}\label{eq:far-field}
    u_n^\infty(\widehat x,\theta)
    =\frac{k^2}{4\pi}\int_B e^{-ik\widehat x\cdot y}(n(y)-1)u_{n,\theta}(y)\,dy,
    \qquad (\widehat x,\theta)\in S^2\times S^2.
\end{equation}
The inverse medium scattering problem involves recovering $n$ in $B$ from the far-field pattern $u_n^\infty$ on $S^2\times S^2$.

This inverse problem is known to be uniquely solvable when we have access to the full far-field pattern on $S^2\times S^2$ \cite{novikov1988multidimensional, nachman1988reconstructions,ramm1988recovery,colton1998inverse,kirsch2011introduction}. It has been shown also that, under certain assumptions, probing the medium with a single planar wave and measuring the corresponding far-field pattern on $S^2$ uniquely determines the values of a perturbation to the refractive
index on the corners of its support \cite{blaasten2020recovering}, and that a polyhedral inclusion can be reconstructed with a single far-field pattern measurement \cite{blaasten2021corners}. We refer also to \cite{di2026discretization} for the inverse problem of reconstructing inhomogeneities by performing a finite number of scattering measurements of acoustic type in the time-harmonic setting. The infinite-dimensional problem has been shown to be only logarithmically stable \cite{stefanov1990stability}. We refer to \cite{bourgeois2013remark} for Lipschitz stability estimates when the unknown is finite-dimensional and the full far-field pattern is available, and to \cite{alberti2022infinite} for a Lipschitz stability estimate for finite-dimensional unknowns with only finitely many measurements. 

We can apply  \cref{thm:main}, in the setting of \Cref{sec:finite-analytic-sampling}, with
\begin{itemize}
    \item $X=L^\infty(B;\mathbb{C})$,
    \item $A=\mathcal A_k$,
    \item $Z=S^2\times S^2$,
    \item ${\rm vol}_Z=m_{S^2}\otimes m_{S^2}$, where $m_{S^2}$ is the surface measure of $S^2$,
    \item and $\mathfrak m\colon \mathcal{A}_k\times (S^2\times S^2)\to \mathbb K$ given by $\mathfrak m(n,(\widehat x,\theta))=u^\infty_n(\widehat x,\theta)$.
\end{itemize}
This leads to a recovery guarantee for finite-dimensional refractive indices using only finitely many incident directions and for each of them measuring the corresponding far-field pattern only at a single location. 

The assumption we need is the injectivity of the forward model on $\mathcal A_k$ with full-data. We notice that the Helmholtz volume potential $V_n$ \eqref{eq:volume-potential} maps $L^2(B)$ continuously into $H^2(B)$ \cite[Theorem 8.2]{colton1998inverse}, hence $V_n\colon L^2(B)\to L^2(B)$ is compact by the Rellich theorem. Since $V_n$ is compact, injectivity of $I-V_n$ implies invertibility by the Fredholm alternative. In particular, every standard admissible class for which the homogeneous radiating problem is uniquely solvable is contained in $\mathcal A_k$; see \cite{colton1998inverse,kirsch2011introduction}. Classical examples are real-valued indices \cite{ramm1988recovery, novikov1988multidimensional} or complex-valued indices with positive imaginary part \cite{bourgeois2013remark}.

\begin{theorem}\label{thm:scattering}
Let $\calM\subset L^\infty(B;\CC)$ be a second-countable embedded real-analytic manifold of dimension $d$. Assume that
\[
    n\longmapsto u_{n}^\infty\in L^2(S^2\times S^2)
\]
is injective on $\mathcal A_k\cap\calM$. Let $\mu$ be a probability measure on $S^2\times S^2$ that is absolutely continuous with respect to the product surface measure. Let $M\geq2d+1$ and let $(\widehat x_1,\theta_1),\dots,(\widehat x_M,\theta_M)\in S^2\times S^2$ be independent samples with law $\mu$. Then, with probability one, for every $n_1,n_2\in\mathcal A_k\cap\calM$,
\[ \text{if }
    u_{n_1}^\infty(\widehat x_i,\theta_i)=u_{n_2}^\infty(\widehat x_i,\theta_i), \; \text{for } i=1,\dots,M, \quad \text{then } n_1=n_2.
\]
\end{theorem}


\section{Proof of the abstract results}\label{sec:proof-main}

\subsection{Analytic and subanalytic preliminaries}
We use the standard notion of dimension of a semianalytic or subanalytic set, defined as the largest dimension of a stratum in an analytic stratification \cite[Remark 2.12]{bierstone1988semianalytic}. In particular, since on smooth manifolds this agrees with Hausdorff dimension, this notion of dimension coincides with the Hausdorff one. We refer to \cite{bierstone1988semianalytic,narasimhan1985analysis} for the relevant background. Finally, we record for later use that the measurement map is analytic in charts: if $\varphi\colon V\subset\RR^d\to X$ is a chart of $\calM$ with $\varphi(V)\subset A$, then $(t,z)\mapsto\mathfrak m\bigl(\varphi(t),z\bigr)$ is real analytic, being a composition of real-analytic maps in the sense of \cref{def:analytic}. This is what allows us to treat the zero sets below as semianalytic subsets of finite-dimensional real-analytic manifolds.

In what follows, we provide a collection of auxiliary results needed for the proofs of our main theorems. Some of these are classical results, reformulated to fit the specific setting of this work.

\begin{lemma}\label[lemma]{lem:analytic-zero-set}
Let $\mathcal N$ be a connected real-analytic manifold of dimension $q$, and let $f\colon\mathcal N\to\KK$ be real-analytic and not identically zero. Then $f^{-1}(0)$ is semianalytic and
\[
    \dim f^{-1}(0)\leq q-1.
\]
\end{lemma}

\begin{proof}
For $\KK=\CC$, replace $f$ by the real-valued analytic function $|f|^2$. The zero level set is semianalytic by definition. In particular, it admits a stratification into analytic submanifolds \cite[Proposition 2.10]{bierstone1988semianalytic}. If the zero set had dimension $q$, then the analytic stratification would contain a $q$-dimensional stratum, and hence the zero set would contain a nonempty open subset of $\mathcal N$. The real-analytic identity principle would then imply $f\equiv0$ on the connected manifold $\mathcal N$, a contradiction.
\end{proof}

\begin{lemma}\label[lemma]{lem:fiber-dimension}
Let $\calX$ and $\calY$ be finite-dimensional real-analytic manifolds. Let $P\subset\calX$ be a relatively compact semianalytic set of dimension $p$, and let
\[
    \widetilde E\subset P\times\calY
\]
be semianalytic. For $x\in P$, write
\[
    E_x=\{y\in\calY:(x,y)\in\widetilde E\}.
\]
If $\dim E_x\leq k$ for every $x\in P$, then
\[
    \dim\widetilde E\leq p+k.
\]
Moreover, denoting by $\pi_\calY\colon\calX\times\calY\to\calY$ the canonical projection, we have that $E=\pi_\calY(\widetilde E)$ is subanalytic and satisfies
\[
    \dim E\leq p+k.
\]
\end{lemma}
\begin{proof}
Choose a locally finite analytic stratification of $\widetilde E$ into analytic
submanifolds \cite{bierstone1988semianalytic}, and let $S$ be one of its strata.
Denote by $\pi_\calX\colon\calX\times\calY\to\calX$ the canonical projection and set
\[
    r=\max_{w\in S}\operatorname{rank}\,d\bigl(\restr{\pi_\calX}{S}\bigr)_w .
\]
Since $\pi_\calX(S)\subset P$ and $\dim P=p$, we have $r\leq p$: at a point where the
rank equals $r$, the image of $\restr{\pi_\calX}{S}$ contains an $r$-dimensional
submanifold of $\calX$ contained in $P$, and the dimension of a subanalytic set is
monotone with respect to inclusion.
 
The set $U\subset S$ on which the rank equals $r$ is nonempty and open in $S$: some
$r\times r$ minor of $d(\restr{\pi_\calX}{S})$ is nonzero at a point where the maximum
is attained, and it remains nonzero in a neighbourhood by continuity. On $U$ the map
$\restr{\pi_\calX}{U}$ has constant rank $r$, so by the constant rank theorem
\cite[\S2.2]{narasimhan1985analysis} each of its fibres is locally an analytic
submanifold of dimension $\dim S-r$. For $x\in P$, the fibre
$(\restr{\pi_\calX}{U})^{-1}(x)$ is contained in $\{x\}\times E_x$, whence
\[
    \dim S-r\ \leq\ \dim(\{x\}\times E_x)=\dim E_x\ \leq\ k .
\]
Together with $r\leq p$ this gives $\dim S\leq p+k$. As $\dim\widetilde E$ is the
maximum of $\dim S$ over the strata, $\dim\widetilde E\leq p+k$.
 
For the second assertion, let $y\in\calY$ and let $V$ be a relatively compact
neighbourhood of $y$, which exists by local compactness of $\calY$. Then
$\widetilde E\cap(P\times V)$ is relatively compact, because $P$ is, and semianalytic;
hence $E\cap V=\pi_\calY\bigl(\widetilde E\cap(P\times V)\bigr)$ is subanalytic by
definition. Since subanalyticity is a local property, $E$ is subanalytic. Finally, the
image of a subanalytic set under a projection has dimension at most that of the set
\cite{bierstone1988semianalytic}, so
\[
    \dim E\ \leq\ \dim\widetilde E\ \leq\ p+k. \qedhere
\]
\end{proof}


\begin{lemma}\label[lemma]{lem:semianalytic-exhaustion}
Every second-countable finite-dimensional real-analytic manifold is a countable union of compact semianalytic subsets.
\end{lemma}

\begin{proof}
Choose a countable analytic atlas and, inside each chart, a countable family of closed Euclidean balls whose interiors cover the chart domain and whose closures remain in that domain. Their inverse images are compact semianalytic sets and form the required cover.
\end{proof}

\subsection{A uniform finite-dimensional truncation}
The following lemma is used only for Gaussian multilinear sampling.

\begin{lemma}\label[lemma]{lem:multilinear-truncation}
Let $Z$ be a separable Hilbert space with orthonormal basis $(e_j)_{j\geq1}$, let $P$ be a compact metric space, and let
\[
    \Phi\colon P\times Z^r\longrightarrow\KK
\]
be continuous. Assume that $\Phi(x,\cdot)$ is a bounded multilinear form and is not identically zero for every $x\in P$. Set $Z_N=\operatorname{span}\{e_1,\ldots,e_N\}$. Then there exists $N$ such that
\[
    \restr{\Phi(x,\cdot)}{Z_N^r}\not\equiv0
    \qquad\text{for every }x\in P.
\]
\end{lemma}

\begin{proof}
Fix $x\in P$. Choose $(z^1,\ldots,z^r)\in Z^r$ with $\Phi(x,z^1,\ldots,z^r)\neq0$. Approximate each $z^\ell$ by a finite linear combination of basis vectors. By continuity, one obtains finite linear combinations $(v^1,\ldots,v^r)$ such that $\Phi(x,v^1,\ldots,v^r)\neq0$. Expanding by multilinearity shows that
\[
    \Phi\bigl(x,e_{j_1(x)},\ldots,e_{j_r(x)}\bigr)\neq0
\]
for at least one tuple of basis vectors occurring in these combinations. By continuity, the same inequality holds in a neighborhood of $x$. A finite subcover of $P$ yields finitely many index tuples; taking $N$ larger than all indices proves the claim.
\end{proof}

\subsection{Proof of the uniform identifiability theorem}
\begin{proof}[Proof of \cref{thm:main}]
Set
\[
    \Delta=\{(a,a):a\in A\cap\calM\},
    \qquad
    \mathcal P
    =\bigl((A\cap\calM)\times(A\cap\calM)\bigr)\setminus\Delta.
\]
This is a second-countable real-analytic manifold of dimension $2d$. By \cref{lem:semianalytic-exhaustion}, choose compact semianalytic sets $K_j\subset\mathcal P$ such that $$\mathcal P=\bigcup_{j\geq1}K_j.$$ By subadditivity of the measure $\nu$, it suffices to show that, for each $j$, the set of measurement systems failing to distinguish a pair in $K_j$ has measure zero.

Fix one such compact set $K=K_j$ and define
\[
    \Phi((a_1,a_2),z)
    =\mathfrak m(a_1,z)-\mathfrak m(a_2,z).
\]
By the relevant separation assumption, $\Phi(p,\cdot)$ is not identically zero for every $p=(a_1,a_2)\in K$. The set of measurement systems failing to distinguish a pair in $K$ is given by
\begin{equation}\label{eq:Ek}
    E_K=\{(z_1,\dots,z_M): \Phi(p,z_i)=0,\,i=1,\dots,M,~\text{for some}~p\in K\}.
\end{equation}

\smallskip
\noindent\emph{Finite-dimensional analytic sampling.}
Let $q=\dim Z$. For $p\in K$, \cref{lem:analytic-zero-set} gives
\[
    \dim\Phi(p,\cdot)^{-1}(0)\leq q-1.
\]
For $M$ measurements, consider the semianalytic incidence set
\[
    \widetilde E_K
    =\left\{(p,z_1,\ldots,z_M)\in K\times Z^M:
      \Phi(p,z_i)=0\ \text{for every }i\right\}.
\]
Notice that $E_K$ is the projection onto $Z^M$ of $\widetilde E_K$. Its fiber over $p$, $\prod_{i=1}^M\Phi(p,\cdot)^{-1}(0)$, has dimension at most $M(q-1)$, and \cref{lem:fiber-dimension} yields
\[
    \dim E_K\leq \dim\widetilde E_K
    \leq 2d+M(q-1).
\]
If $M\geq2d+1$, then
\[
    \dim E_K\leq qM-1.
\]
Thus $E_K$ has zero $qM$-dimensional volume. Since $\mu^M$ is absolutely continuous with respect to the product volume measure, $\mu^M(E_K)=0$.

\smallskip
\noindent\emph{Gaussian multilinear sampling.}
Let $Q$ be the covariance operator of $\gamma$. Since $Q$ is positive, trace class, and injective, there is an orthonormal eigenbasis $(e_j)_{j\geq1}$ with strictly positive eigenvalues. Apply \cref{lem:multilinear-truncation} to $\Phi$ on $K$ and obtain $N$ such that the restriction of every polynomial $\Phi(p,\cdot)$ to $Z_N^r$ is nonzero.

Let $P_N$ denote the orthogonal projection of $Z$ onto $Z_N$ and set
\[
    \gamma_N=(P_N)_\#\gamma,
    \qquad
    \gamma_N^\perp=(I-P_N)_\#\gamma .
\]
Because $Z_N$ is spanned by eigenvectors of $Q$, both $Z_N$ and $Z_N^\perp$ are
$Q$-invariant, and $\langle Qu,v\rangle=0$ whenever $u\in Z_N$ and $v\in Z_N^\perp$. The
jointly Gaussian real random variables $z\mapsto\langle z,u\rangle$ and
$z\mapsto\langle z,v\rangle$ are then uncorrelated, hence independent; consequently $P_Nz$
and $(I-P_N)z$ are independent under $\gamma$, and, under the identification
$Z\cong Z_N\times Z_N^\perp$,
\begin{equation}\label{eq:gaussian-splitting}
    \gamma=\gamma_N\otimes\gamma_N^\perp .
\end{equation}
Moreover $\gamma_N$ is a Gaussian measure on the $N$-dimensional space $Z_N$ whose
covariance $\restr{Q}{Z_N}$ has the strictly positive eigenvalues
$\lambda_1,\ldots,\lambda_N$; it is therefore nondegenerate and has a smooth positive
density with respect to Lebesgue measure on $Z_N$. Note that the splitting
\eqref{eq:gaussian-splitting} uses that $(e_j)_{j\geq1}$ diagonalises $Q$: for an arbitrary
orthonormal basis the two components need not be independent. The mean of $\gamma$ plays no
role in the argument.

The bad set $E_K\subset(Z^r)^M$ \eqref{eq:Ek} is closed: this follows from compactness of $K$ and continuity of $\Phi$.

Fix tails $w_i\in(Z_N^\perp)^r$, $i=1,\ldots,M$. For $p=(a_1,a_2)\in K$ and $i=1,\dots,M$, the function
\[
    z\longmapsto\Phi(p,z+w_i),\qquad z\in Z_N^r,
\]
can be written as
\begin{align*}
    \Phi(p,z+w_i)
    &=\mathfrak{m}(a_1,z+w_i)-\mathfrak{m}(a_2,z+w_i)\\
\end{align*}
where we used the definition of $\mathfrak m$ \eqref{eq:gaussian-measurement-map}. By exploiting the multilinearity of $\mathfrak m(a_1,\cdot)-\mathfrak m(a_2,\cdot)$, we can write this expression compactly. For every $J\subseteq\{1,\dots,r\}$, let $s_J\in H^r$ defined as
\[
    s_J^j = 
    \begin{cases} 
         z^j  & \text{if } j \in J, \\ 
         w_i^j & \text{if } j \notin J. 
    \end{cases}
\]
The full expression reads as follows
\[  
    \Phi(p,z+w_i)=\sum_{J\subseteq\{1,\dots,r\}}\mathfrak m(a_1,s_J)-\mathfrak m(a_2,s_J).
\]

We split it into two parts, separating the term with $J=\{1,\dots,r\}$, and the remaining terms, namely we define
\[
    P(z)=\mathfrak{m}(a_1,z)-\mathfrak{m}(a_2,z),
\]
and
\[
    Q(z)= \sum_{J\subsetneq\{1,\dots,r\}}\mathfrak m(a_1,s_J)-\mathfrak m(a_2,s_J),
\]
and we have $\Phi(p,z+w_i)=P(z)+Q(z)$, for $z\in Z_{N}^r$.
Notice that, by the choice of $N$, $P$ is a homogeneous polynomial of degree $r$ that is not identically zero on $Z_{N}^r$. On the other hand, $Q$ is a polynomial of degree strictly smaller than $r$. Indeed, in every term of the sum defining $Q$, there is at least one fixed entry in $Z_{N}^\perp$. In particular, homogeneous polynomials of different degrees are linearly independent as elements of the polynomial ring. Consequently, their sum $\Phi$ cannot vanish identically on $Z_{N}^r$.

Consequently, each of its zero sets has dimension at most $Nr-1$. The sliced incidence set
\[
    \widetilde E_K(w_1,\ldots,w_M)
    =\left\{
      \begin{aligned}
        &(p,z_1,\ldots,z_M)\in K\times(Z_N^r)^M:\\[-2pt]
        &\Phi(p,z_i+w_i)=0,~i=1,\dots,M
      \end{aligned}
    \right\}
\]
therefore has dimension at most
\[
    2d+M(Nr-1)\leq MNr-1,
\]
where we used $M\ge2d+1$.
Its projection $E_K(w_1,\ldots,w_M)$ onto $(Z_N^r)^M$ is subanalytic of the same or smaller dimension. Since $(\gamma_N^r)^M$ has a smooth density with respect to Lebesgue measure,
\[
    (\gamma_N^r)^M\bigl(E_K(w_1,\ldots,w_M)\bigr)=0.
\]
Fubini's theorem then gives 
\[
   (\gamma^r)^M(E_K)=\int_{(Z_N^\perp)^{rM}} (\gamma_N^r)^M\bigl(E_K(w_1,\ldots,w_M)\bigr) d(\gamma_N^\perp)^{rM}(w_1,\dots,w_M)=0.
\]

Finally, the full bad set is contained in the countable union $\bigcup_{j\geq1}E_{K_j}$, and therefore has probability zero.
\end{proof}

\begin{proof}[Proof of \cref{cor:nonuniform}]
Repeat the proof with the parameter set
\[
    (A\cap\calM)\setminus\{a^\dagger\}
\]
in place of the pair manifold $\mathcal P$. Its dimension is $d$, so the same calculation gives a null bad set as soon as $M>d$.
\end{proof}

\begin{proof}[Proof of \cref{thm:sparse}]
If $M\geq2d+1$, the conclusion follows from \cref{thm:main}. Suppose instead that $M\geq4s+1$. For any $a_1,a_2\in\calM_s$, the union of their supports is contained in a set $J\subset\{1,\ldots,d\}$ with $|J|\leq2s$. Both parameters lie in
\[
    \calM_J=\operatorname{span}\{e_j:j\in J\},
\]
whose dimension is at most $2s$. By \cref{thm:main}, the bad measurement systems for $A\cap\calM_J$ form a null set. There are only finitely many possible sets $J$, and the bad set for $\calM_s$ is contained in the finite union of these null sets. Combining the two regimes gives \eqref{eq:sparse-bound}.
\end{proof}

\section{The Calder\'on forward map and finite-dimensional truncation}\label{sec:proof-calderon}

\subsection{Analyticity and compactness}
Let
\[
    \mathcal H_\diamond
    =\left\{u\in H^1(\Omega):\int_{\partial\Omega}\restr{u}{\partial\Omega}\,ds=0\right\}.
\]
For $\sigma\in L^\infty_+(\Omega)$, define
\[
    A_\sigma\colon\mathcal H_\diamond\to\mathcal H_\diamond^*,
    \qquad
    (A_\sigma u)(v)=\int_\Omega\sigma\nabla u\cdot\nabla v\,dx.
\]
Let $J\colon L^2_\diamond(\partial\Omega)\to\mathcal H_\diamond^*$ be given by
\[
    (Jf)(v)=\int_{\partial\Omega}f\,\restr{v}{\partial\Omega}\,ds,
\]
and let $T\colon\mathcal H_\diamond\to L^2_\diamond(\partial\Omega)$ be the trace map.

\begin{proposition}\label[proposition]{prop:nd-analytic}
The map
\[
    \Lambda\colon L^\infty_+(\Omega)\longrightarrow
    \calK\bigl(L^2_\diamond(\partial\Omega)\bigr),
    \qquad \sigma\longmapsto\Lambda_\sigma,
\]
is real-analytic. Its differential is
\begin{equation}\label{eq:nd-derivative-operator}
    d\Lambda_\sigma(\eta)
    =-T A_\sigma^{-1}A_\eta A_\sigma^{-1}J,
\end{equation}
and, equivalently,
\begin{equation}\label{eq:nd-derivative-form}
    \bigl\langle d\Lambda_\sigma(\eta)f,g\bigr\rangle
    =-\int_\Omega \eta\,\nabla u_{\sigma,f}\cdot\nabla u_{\sigma,g}\,dx.
\end{equation}
Here $A_\eta$ is defined by the same bilinear form as $A_\sigma$. In particular, $d\Lambda_\sigma(\eta)$ is compact for every $\sigma\in L^\infty_+(\Omega)$ and $\eta\in L^\infty(\Omega)$.
\end{proposition}

\begin{proof}
The Poincar{\'e} inequality on $\mathcal H_\diamond$, coercivity, and the Lax--Milgram theorem give $\Lambda_\sigma=TA_\sigma^{-1}J$. The map $\sigma\mapsto A_\sigma$ is affine and continuous, and inversion is analytic on the open set of invertible bounded operators. The trace maps continuously into $H^{1/2}(\partial\Omega)$, and the embedding $H^{1/2}(\partial\Omega)\hookrightarrow L^2(\partial\Omega)$ is compact; hence $T$ is compact. The bounded linear map $S\mapsto TSJ$ takes bounded operators $\mathcal H_\diamond^*\to\mathcal H_\diamond$ into compact operators on $L^2_\diamond(\partial\Omega)$. Consequently, $\Lambda_\sigma$ is compact and $\sigma\mapsto\Lambda_\sigma$ is analytic with values in the operator norm. Differentiating the inverse gives \eqref{eq:nd-derivative-operator}, and the variational formulation yields \eqref{eq:nd-derivative-form}.
\end{proof}

\begin{proof}[Proof of \cref{thm:calderon-gaussian}]
Apply the Gaussian architecture with
\[
    H=Z=L^2_\diamond(\partial\Omega),\qquad r=2,\qquad\psi=I_{L^2_\diamond(\partial\Omega)},
\]
and with $F(\sigma)$ equal to the bounded bilinear form 
\begin{equation*}
    \Lambda_\sigma(f,g)=\langle \Lambda_\sigma f,g\rangle_{L^2(\partial\Omega)}.
\end{equation*}
Analyticity follows from \cref{prop:nd-analytic}, and injectivity is assumed. The conclusion is \cref{thm:main}.
\end{proof}

\subsection{Uniform preservation of injectivity under truncation}\label{subsection:uniform-preservation-injectivity-truncation}
The following argument is modeled on the compactness mechanisms used in \cite{alberti2022infinite,alberti2022inverse}, but it is stated here in a form that simultaneously preserves injectivity of the forward map and injectivity of its differential.

\begin{lemma}\label[lemma]{lem:compact-secants}
Let $\calM\subset X$ be a finite-dimensional $C^1$ embedded manifold, let $U\subset\calM$ be open, let $K\subset U$ be compact, and let $G\in C^1(U;Y)$, where $Y$ is a Banach space. Then the normalized secant set
\[
    \calS_K(G)
    =\left\{\frac{G(x)-G(y)}{\norm{x-y}_X}:x,y\in K,\ x\neq y\right\}
\]
is relatively compact in $Y$. Moreover,
\[
    \calT_K(G)
    =\left\{dG_x(v):x\in K,\ v\in T_x\calM,\ \norm{v}_X=1\right\}
\]
is compact.

If $G$ is injective on $K$ and $dG_x$ is injective for every $x\in K$, then
\begin{equation}\label{eq:secant-lower-bound}
    0\notin\overline{\calS_K(G)}\cup\calT_K(G).
\end{equation}
\end{lemma}

\begin{proof}
Take a sequence of normalized secants $\frac{G(x_j)-G(y_j)}{\norm{x_j-y_j}_X}$ for $x_j,y_j\in K$. After passing to a subsequence, the $x_j,y_j$ converge to $x,y\in K$ respectively. If $x\neq y$, convergence is immediate. If $x=y$, choose a $C^1$ chart $\varphi:V\subset\RR^d\to\calM$ around $x$, with $V$ convex after shrinking. Write the endpoints as $x_j=\varphi(p_j)$ and $y_j=\varphi(q_j)$, where $p_j,q_j\to q=\varphi^{-1}(x)$, and pass to a subsequence such that
\[
    \frac{p_j-q_j}{|p_j-q_j|}\longrightarrow v,\qquad |v|=1.
\]
The fundamental theorem of calculus gives
\[
    \frac{G(\varphi(p_j))-G(\varphi(q_j))}{|p_j-q_j|}
    \longrightarrow d(G\circ\varphi)_{q}(v),
\]
and
\[
    \frac{\norm{\varphi(p_j)-\varphi(q_j)}_X}{|p_j-q_j|}
    \longrightarrow \norm{d\varphi_q(v)}_X>0.
\]
Consequently,
\[
    \frac{G(\varphi(p_j))-G(\varphi(q_j))}
         {\norm{\varphi(p_j)-\varphi(q_j)}_X}
    \longrightarrow
    \frac{d(G\circ\varphi)_q(v)}{\norm{d\varphi_q(v)}_X},
\]
so the normalized secants have a convergent subsequence. Hence, $\mathcal S_K(G)$ is relatively compact in $Y$. The unit tangent bundle over $K$, with the norm induced by $X$, is compact by local triviality and a finite chart cover of $K$; continuity of $dG$ therefore gives compactness of $\calT_K(G)$.

If a limit of secants with distinct limiting endpoints were zero, injectivity of $G$ on $K$ would fail. If the endpoints coalesce, the preceding calculation shows that the limit is a nonzero normalized derivative whenever $dG_x$ is injective. The same injectivity excludes zero from $\calT_K(G)$.
\end{proof}

\begin{lemma}\label[lemma]{lem:uniform-on-compact}
Let $Y_0$ be a Banach space and let $Q_N\in\calL(Y_0)$ satisfy
\[
    \sup_N\norm{Q_N}<\infty,
    \qquad Q_Ny\longrightarrow y\quad\text{for every }y\in Y_0.
\]
Then $Q_N\to I$ uniformly on every compact subset of $Y_0$.
\end{lemma}

\begin{proof}
Let $C\subset Y_0$ be compact and set $B=\sup_N\norm{I-Q_N}<\infty$. Fix $\varepsilon>0$. If $B=0$, there is nothing to prove. Otherwise, choose a finite set $\{y_1,\ldots,y_J\}\subset C$ such that every $y\in C$ satisfies
\[
    \norm{y-y_j}<\frac{\varepsilon}{2B}
\]
for some $j$. Pointwise convergence gives $N_0$ such that
\[
    \norm{(I-Q_N)y_j}<\frac{\varepsilon}{2},
    \qquad j=1,\ldots,J,\ N\geq N_0.
\]
For such $N$ and $y\in C$, choose $j$ as above. Then
\[
    \norm{(I-Q_N)y}
    \leq B\norm{y-y_j}+\norm{(I-Q_N)y_j}
    <\varepsilon.
\]
Thus the convergence is uniform on $C$.
\end{proof}

\begin{proposition}\label[proposition]{prop:truncation-preserves-injectivity}
In the setting of \cref{lem:compact-secants}, assume that $G(U)\subset Y_0$, where $Y_0$ is a closed subspace of $Y$, and let $Q_N\in\calL(Y_0)$ be uniformly bounded with $Q_Ny\to y$ in $Y_0$ for every $y\in Y_0$. Since $Y_0$ is closed, one has $dG_x(T_x\calM)\subset Y_0$ for every $x\in U$. If $G$ is injective on $K$ and $dG_x$ is injective for every $x\in K$, then, for all sufficiently large $N$,
\begin{enumerate}[label=\textup{(\roman*)},leftmargin=*]
    \item $Q_NG$ is injective on $K$;
    \item $Q_NdG_x:T_x\calM\to Y_0$ is injective for every $x\in K$.
\end{enumerate}
More precisely, there is $c>0$ such that, for all sufficiently large $N$,
\[
    \norm{Q_NG(x)-Q_NG(y)}_Y\geq c\norm{x-y}_X,
    \qquad x,y\in K.
\]
\end{proposition}

\begin{proof}
By \cref{lem:compact-secants} and \eqref{eq:secant-lower-bound}, the set
\[
    C=\overline{\calS_K(G)}\cup\calT_K(G)
\]
is compact and does not contain zero. If $C=\varnothing$, the conclusions are immediate. Otherwise, set
\[
    \delta=\min_{v\in C}\norm{v}_Y>0.
\]
By \cref{lem:uniform-on-compact}, $Q_N$ converges uniformly to the identity on $C$. Hence, for all sufficiently large $N$,
\[
    \sup_{v\in C}\norm{Q_Nv-v}_Y<\frac{\delta}{2},
\]
and therefore $\norm{Q_Nv}_Y\geq\frac12\norm{v}_Y$ for every $v\in C$. Applying this estimate to normalized secants gives the asserted lower bound with $c=\delta/2$, and applying it to normalized tangent vectors gives injectivity of every $Q_NdG_x$.
\end{proof}

\begin{lemma}\label[lemma]{lem:injective-extension}
Let $\calM$ be a finite-dimensional $C^1$ manifold, let $Y$ be a Banach space, and let $G\in C^1(\calM;Y)$. If $K\subset\calM$ is compact, $G$ is injective on $K$, and $dG_x$ is injective for every $x\in K$, then there is an open neighborhood $U$ of $K$ in $\calM$ such that $G$ is injective on $U$.
\end{lemma}

\begin{proof}
Fix $x\in K$ and a local chart $\varphi$ at $x$. Since $d(G\circ\varphi)$ is injective at the chart point $x$, its finite-dimensional range is complemented in $Y$. Thus there is a bounded linear map $L\colon Y\to\RR^d$ such that $L\circ d(G\circ\varphi)$ is invertible there. The inverse function theorem applied to $L\circ G\circ\varphi$ shows that $G$ is injective on a neighborhood of $x$.

Choose a metric compatible with the topology of $\calM$. If no neighborhood of $K$ had the asserted property, there would be distinct $x_j,y_j$ with $\dist(x_j,K)+\dist(y_j,K)\to0$ and $G(x_j)=G(y_j)$. Compactness gives subsequences converging to $x,y\in K$. Continuity and injectivity on $K$ imply $x=y$. For large $j$, both points then lie in a neighborhood on which $G$ is locally injective, a contradiction.
\end{proof}

\subsection{Proof of the finite-dimensional Calder\'on theorem}
\begin{proof}[Proof of \cref{thm:calderon-finite}]
Let
\[
    \calK_2=\calK\bigl(L^2_\diamond(\partial\Omega)\bigr)
\]
and define
\[
    Q_NT=P_NTP_N,
    \qquad T\in\calK_2,
\]
where $P_N$ is the orthogonal projection from $L^2_\diamond(\partial\Omega)$ to $H_N$ \eqref{eq HN}.
Then $\norm{Q_N}\leq1$ and
\begin{equation}\label{eq:compact-operator-truncation}
    \norm{T-P_NTP_N}\longrightarrow0
    \qquad\text{for every }T\in\calK_2.
\end{equation}
Indeed,
\[
    \norm{T-P_NTP_N}
    \leq\norm{(I-P_N)T}+\norm{T(I-P_N)},
\]
and both terms tend to zero because $T$ and $T^*$ are compact and $P_N\to I$ strongly.

By \cref{prop:nd-analytic}, the forward map and all its differentials take values in $\calK_2$. Apply \cref{prop:truncation-preserves-injectivity} to
\[
    G(\sigma)=\Lambda_\sigma,
    \qquad \sigma\in L^\infty_+(\Omega)\cap\calM,
\]
and to the operators $Q_N$. For every sufficiently large $N$, the truncated map
\[
    G_N(\sigma)=P_N\Lambda_\sigma P_N
\]
is injective on $K$, and $dG_{N,\sigma}=Q_Nd\Lambda_\sigma$ is injective on $T_\sigma\calM$ for every $\sigma\in K$. By \cref{lem:injective-extension}, $G_N$ is injective on an open neighborhood $U_N$ of $K$ in $L^\infty_+(\Omega)\cap\calM$. Since $\calM$ carries the subspace topology and $L^\infty_+(\Omega)$ is open in $L^\infty(\Omega)$, there is an open set $A_N\subset L^\infty(\Omega)$, with $A_N\subset L^\infty_+(\Omega)$, such that $U_N=A_N\cap\calM$.

The scalar map
\[
    (\sigma,f,g)\longmapsto\langle\Lambda_\sigma f,g\rangle,
    \qquad (\sigma,f,g)\in A_N\times H_N\times H_N,
\]
is real-analytic. Equality of these scalar values for all $(f,g)\in H_N\times H_N$ is equivalent to equality of $G_N(\sigma)$. Hence the separation condition \eqref{eq:analytic-separation} holds on $U_N=A_N\cap\calM$. Apply \cref{thm:main} with sampling manifold $H_N\times H_N$ and probability $\rho_N\otimes\rho_N$. Since $K\subset U_N$, the resulting random measurements are injective on $K$ almost surely.
\end{proof}

\section{Proof of the scattering result}\label{sec:proof-scattering}

\begin{proposition}\label[proposition]{prop:far-field-analytic}
The map
\[
    \mathcal A_k\times S^2\times S^2\longrightarrow\CC,
    \qquad
    (n,\widehat x,\theta)\longmapsto u_n^\infty(\widehat x,\theta),
\]
is real-analytic.
\end{proposition}

\begin{proof}
The map $n\mapsto V_n$ in \eqref{eq:volume-potential} is affine and continuous from $L^\infty(B;\CC)$ into $\calL(L^2(B))$. The set of invertible bounded operators is open, and inversion is analytic there. Therefore
\[
    n\longmapsto (I-V_n)^{-1}
\]
is analytic on $\mathcal A_k$. The map $\theta\mapsto e^{ik(\cdot)\cdot\theta}$ is real-analytic from $S^2$ into $L^2(B)$, so \eqref{eq:lippmann-schwinger} shows that $(n,\theta)\mapsto u_{n,\theta}$ is jointly real-analytic as an $L^2(B)$-valued map.

Similarly, $\widehat x\mapsto e^{-ik\widehat x\cdot(\cdot)}$ is real-analytic as an $L^2(B)$-valued map. Formula \eqref{eq:far-field} is obtained by composing these maps with the continuous trilinear map
\[
    L^2(B)\times L^\infty(B)\times L^2(B)\to\CC,\qquad(\phi,m,v)\longmapsto\int_B\phi(y)m(y)v(y)\,dy.
\]
The conclusion follows.
\end{proof}

\begin{proof}[Proof of \cref{thm:scattering}]
Use the finite-dimensional analytic sampling architecture with
\[
    Z=S^2\times S^2,
    \qquad
    \mathfrak m(n,(\widehat x,\theta))=u_n^\infty(\widehat x,\theta).
\]
Joint analyticity follows from \cref{prop:far-field-analytic}. If $n_1\neq n_2$ in $\mathcal A_k\cap\calM$, injectivity of the full far-field map implies that
\[
    (\widehat x,\theta)\longmapsto
    u_{n_1}^\infty(\widehat x,\theta)-u_{n_2}^\infty(\widehat x,\theta)
\]
is not identically zero. Thus \eqref{eq:analytic-separation} holds, and \cref{thm:main} gives the result.
\end{proof}

\section*{Acknowledgements}
We thank Giacomo Traverso for his contributions in the early phase of the project. The research was supported in part by the MIUR Excellence Department Project awarded to Dipartimento di Matematica, Università di Genova, CUP D33C23001110001. Co-funded by the European Union (ERC, SAMPDE, 101041040). Views and opinions expressed are however those of the authors only and do not necessarily reflect those of the European Union or the European Research Council. Neither the European Union nor the granting authority can be held responsible for them.


\bibliographystyle{abbrvnat}
\bibliography{ref}

\end{document}
